\documentclass[a4paper, 10pt]{article}
%\usepackage{fontspec}
%\setmainfont{Lato}
\usepackage{pgf}
\usepackage{eurosym}
\usepackage{graphicx}
\usepackage{wasysym}
\usepackage{hyperref}
\usepackage{listings}
\usepackage{pxfonts}
\usepackage{verbatim}
\usepackage{color}
\usepackage{xcolor}
\usepackage{wrapfig}
\usepackage{enumitem}
\usepackage{booktabs}
\usepackage{gensymb}
\usepackage{tabularx}
\usepackage{currfile}

\hypersetup{
    bookmarks=true,         % show bookmarks bar?
    unicode=true,          % non-Latin characters in Acrobat’s bookmarks
    pdftoolbar=true,        % show Acrobat’s toolbar?
    pdfmenubar=true,        % show Acrobat’s menu?
    pdffitwindow=true,     % window fit to page when opened
    pdftitle={Assessments},    % title
    pdfauthor={Paul Vesey},     % author
    pdfsubject={Advanced Graphics Assignment },   % subject of the document
    pdfcreator={},   % creator of the document
    pdfproducer={xelatex}, % producer of the document
    pdfkeywords={'Graphics' }, % list of keywords
    pdfnewwindow=true,      % links in new PDF window
    colorlinks=true,       % false: boxed links; true: colored links
    linkcolor=violet,          % color of internal links (change box color with linkbordercolor)
    citecolor=magenta,        % color of links to bibliography
    filecolor=red,      % color of file links
    urlcolor=blue           % color of external links
}

\setlength\parindent{0pt}
\begin{document}

\lstset{language=HTML,
				basicstyle=\small,
				breaklines=true,
        numbers=left,
        numberstyle=\tiny,
        showstringspaces=false,
        aboveskip=-20pt,
        frame=leftline
        }
				
\begin{table}%
	\begin{minipage}{0.4\textwidth}%
			\includegraphics[width=1\textwidth]{./img/LITlogo.jpg}
	\end{minipage}
	\qquad
	\centering
	\parbox{0.4\textwidth}{
		\begin{large}			
			\begin{tabular}{| r | l |} \hline
				Subject: & \textbf{Advanced Graphics}\\
								 & \textbf{\& Visualisation}\\
				Course: & \textbf{Interior Design Y3}\\
				Session: & \textbf{Autumn 2020}\\
				Lecturer: & \textbf{Paul Vesey \footnotesize{BEng, MIE, HDip}}\\
				Filename: & \footnotesize{\currfilename}\\
				\hline
			\end{tabular}
		\end{large}			
	}
\end{table}
\vspace{0.25cm}	
%    _            _                                  _     _ 
%   / \   ___ ___(_) __ _ _ __  _ __ ___   ___ _ __ | |_  / |
%  / _ \ / __/ __| |/ _` | '_ \| '_ ` _ \ / _ \ '_ \| __| | |
% / ___ \\__ \__ \ | (_| | | | | | | | | |  __/ | | | |_  | |
%/_/   \_\___/___/_|\__, |_| |_|_| |_| |_|\___|_| |_|\__| |_|
%                   |___/                                    

\newpage
\setcounter{page}{1}
\begin{center}
	\begin{figure}[ht]
		\centering
		\includegraphics[width = 6cm]{img/LITlogo.jpg}
		\label{fig:logoa1}
	\end{figure}
	\Large\textbf{Practical 1 - Revit Revision}\\
	\large\textbf{BA in Interior Design \& Technology}
\end{center}


In this practical you are required to create a simple space that we will use later in Unity 3D.  The space should consist of 4 rooms and a central corridor.  The room configuration is at your discretion.  Each room should have a small degree of furnishing and a color scheme to easily distinguish one room from another.  You will have to create \textit{inter-alia} the following model components


\begin{enumerate}
	\item Levels
	\item Floor, Walls, Ceilings and/or Roof
	\item Windows and Doors
\end{enumerate}

You do not need to model any external elements, such as toposurfaces, trees etc.  We will be adding these items later using Unity.

\end{document}





